\documentclass{book}

\usepackage{fontspec}
\setmainfont{Ubuntu}

\usepackage[left=1in, right=1in, top=1in, bottom=1in]{geometry}

\usepackage{fancyhdr}
\pagestyle{fancy}
\lhead{\href{https://github.com/QubitPi/general-relativity}{Supplementary Study Notes }}
\chead{General Relativity}
\rhead{
    \href{https://github.com/QubitPi}{QubitPi}
    \includegraphics[scale=0.05]{logo-8th-version.png}
}
\lfoot{}
\cfoot{}
\rfoot{\thepage}
\renewcommand{\headrulewidth}{0.4pt}
\renewcommand{\footrulewidth}{0.4pt}

\usepackage{hyperref}
\hypersetup{
    colorlinks=true,
    linkcolor=blue,
    anchorcolor=blue,
    urlcolor=blue
}

\usepackage{graphicx}
\usepackage{float}
\graphicspath{ {./img/} }
\usepackage{tikz}
\usepackage[most]{tcolorbox}


\newtcbtheorem[number within=section]{mytheorem}{Theorem}
{colback=green!5,colframe=green!35!black,fonttitle=\bfseries}{th}

\newtcbox{\roundinlinebox}[1][red]{on line,
arc=7pt,colback=#1!10!white,colframe=#1!50!white,
before upper={\rule[-3pt]{0pt}{10pt}},boxrule=1pt,
boxsep=0pt,left=6pt,right=6pt,top=2pt,bottom=2pt}

% https://tex.stackexchange.com/a/307436/277953
\newtcolorbox{marker}[1][]{enhanced,
before skip=2mm,after skip=3mm,
boxrule=0.4pt,left=5mm,right=2mm,top=1mm,bottom=1mm,
colback=yellow!50,
colframe=yellow!20!black,
sharp corners,rounded corners=southeast,arc is angular,arc=3mm,
underlay={%
    \path[fill=tcbcolback!80!black] ([yshift=3mm]interior.south east)--++(-0.4,-0.1)--++(0.1,-0.2);
    \path[draw=tcbcolframe,shorten <=-0.05mm,shorten >=-0.05mm] ([yshift=3mm]interior.south east)--++(-0.4,-0.1)--++(0.1,-0.2);
    \path[fill=yellow!50!black,draw=none] (interior.south west) rectangle node[white]{\Huge\bfseries !} ([xshift=4mm]interior.north west);
},
drop fuzzy shadow,#1}

\usepackage{amsfonts}
\usepackage{parskip}

\setlength{\parindent}{0pt}
